\subsection{Recognizing Direct Products}

Let $G$ be a group.

\begin{exercise}[5.4.1]
    Prove that if $x, y \in G$, then $[y, x] = [x, y]\inv$. Deduce that for any subsets $A$ and $B$ of $G$, $[A, B] = [B, A]$ (recall that $[A, B]$ is the \emph{subgroup} of $G$ generated by the commutators $[a, b]$).
\end{exercise}

\begin{sol}
    Note that
    \[[x, y]\inv = (x\inv y\inv xy)\inv = yx y\inv x\inv = [y, x].\]
    Hence, $[y, x] = [x, y]\inv$. For any subsets $A$ and $B$ of $G$, the subgroup $[A, B]$ is generated by the commutators $[a, b]$ for $a \in A$ and $b \in B$. Since $[a, b] = [b, a]\inv$, it follows that $[A, B] = [B, A]$.
\end{sol}

\begin{exercise}[5.4.2]
    Prove that a subgroup $H$ of $G$ is normal if and only if $[G, H] \leq H$.
\end{exercise}

\begin{sol}
    By \exref{5.4.1}, we know $[H, G] = [G, H]$. Then use Proposition 5.7(2).
\end{sol}

\begin{exercise}[5.4.3]
    Let $a, b, c \in G$. Prove that
    \begin{subexercises}
        \item $[a, bc] = [a, c](c\inv [a, b]c)$,
        \item $[ab, c] = (b\inv[a, c]b)[b, c]$.
    \end{subexercises}
\end{exercise}

\begin{solenum}
    \item $[a, bc] = a\inv (bc)\inv a(bc) = a\inv c\inv b\inv a b c = a\inv c\inv a c (c\inv a\inv b\inv a b c) = [a, c](c\inv [a, b]c)$.
    \item $[ab, c] = (ab)\inv c\inv (ab)c = b\inv a\inv c\inv a b c = b\inv a\inv c\inv a c (c\inv a\inv b\inv a b c) = (b\inv[a, c]b)[b, c]$. \qh
\end{solenum}

\begin{exercise}[5.4.4]
    Find the commutator subgroups of $S_4$ and $A_4$.
\end{exercise}

\begin{sol}
    Note that $A_4 \nsub S_4$ and $S_4/A_4 \cong Z_2$, which is abelian, so $[S_4, S_4] \leq A_4$. On the other hand, $A_4$ is generated by the 3-cycles, and $[S_4, S_4]$ contains all 3-cycles since for any 3-cycle $(a\ b\ c)$, we have that
    \[(a\ b\ c) = (a\ b)\inv (a\ c)\inv (a\ b) (a\ c) = [(a\ c), (a\ b)],\]
    so $A_4 \leq [S_4, S_4]$. Thus, $[S_4, S_4] = A_4$. 

    By \exref{3.5.9}, we know that $V_4 \nsub A_4$ and $A_4/V_4 \cong Z_3$, which is abelian, so $[A_4, A_4] \leq V_4$. On the other hand, $A_4$ is non-abelian, so $[A_4, A_4]$ is non-trivial. Since $V_4$ is the only non-trivial subgroup of $A_4$ that is normal, we must have $[A_4, A_4] = V_4$. \qh
\end{sol}

\begin{exercise}[5.4.5]
    Prove that $A_n$ is the commutator subgroup of $S_n$ for all $n \geq 5$.
\end{exercise}

\begin{sol}
    Similar argument as in the previous exercise, where $S_n/A_n \cong Z_2$.
\end{sol}

\newpage

\begin{exercise}[5.4.6]
    Exhibit a representative of each cycle type of $A_5$ as a commutator in $S_5$.
\end{exercise}

\begin{sol}
    The cycle types of $A_5$ are the identity, 3-cycles, products of two disjoint 2-cycles, and 5-cycles. The identity is the commutator of any element with itself, so we can take $[1, 1] = 1$. For a 3-cycle $(a\ b\ c)$, we have that
    \[(a\ b\ c) = (a\ b)\inv (a\ c)\inv (a\ b) (a\ c) = [(a\ c), (a\ b)].\]
    For a product of two disjoint 2-cycles $(a\ b)(c\ d)$, we have that
    \[(a\ b)(c\ d) = (a\ b\ d)\inv (b\ c\ d)\inv (a\ b\ d) (b\ c\ d) = [(a\ b\ d), (b\ c\ d)].\]
    For a 5-cycle $(a\ b\ c\ d\ e)$, we have that
    \[(a\ b\ c\ d\ e) = (a\ e\ d)\inv ((b\ e)(c\ d))\inv (a\ e\ d) ((b\ e)(c\ d)) = [(a\ e\ d), (b\ e)(c\ d)]. \qh\]
\end{sol}

\begin{exercise}[5.4.7]
    Prove that if $p$ is a prime and $P$ is a non-abelian group of order $p^3$, then $P' = Z(P)$.
\end{exercise}

\begin{sol}
    Let $P$ be a non-abelian group of order $p^3$. By the class equation, we have that $|Z(P)|$ divides $p^3$ and is greater than 1, so $|Z(P)| = p$ or $p^2$. If $|Z(P)| = p^2$, then $P/Z(P)$ is cyclic, so $P$ is abelian, which is a contradiction. Thus, $|Z(P)| = p$, so $P/Z(P)$ is of order $p^2$, hence abelian. Therefore, $P' \leq Z(P)$. On the other hand, if $P' < Z(P)$, then properness of $P'$ implies that $|P'| = 1$, so $P$ is abelian, which is a contradiction. Thus, $P' = Z(P)$.
\end{sol}

\begin{exercise}[5.4.8]
    Assume $x, y \in G$, and both $x$ and $y$ commute with $[x, y]$. Prove that for all $n \in \zp$,
    \[(xy)^n = x^n y^n [y, x]^{\frac{n(n-1)}{2}}.\]
\end{exercise}

\begin{sol}
    We first prove that for all $n \in \zp$,
    \[y^nx = xy^n [y, x]^n.\]
    We do this by induction on $n$. For $n = 1$, we have $yx = (xyy\inv x\inv)(yx) = xy[y, x]$. Now, assume that $y^nx = xy^n [y, x]^n$ for some $n \in \zp$. Then
    \[y^{n+1}x = y(y^nx) = y(xy^n [y, x]^n) = (yx)y^n [y, x]^n = xy[y, x]y^n [y, x]^n = xy^{n+1} [y, x]^{n+1}.\]
    Now, we prove the main statement by induction on $n$. For $n = 1$, we have $(xy)^n = xy = x^n y^n [y, x]^{\frac{n(n-1)}{2}}$. Now, assume that $(xy)^n = x^n y^n [y, x]^{\frac{n(n-1)}{2}}$ for some $n \in \zp$. Then
    \begin{align*}
        (xy)^{n+1} & = (xy)^n (xy) \\
        & = (x^n y^n [y, x]^{\frac{n(n-1)}{2}})(xy) \\
        & = x^n y^n [y, x]^{\frac{n(n-1)}{2}} x y \\
        & = x^n y^n x y [y, x]^{\frac{n(n-1)}{2}} \\
        & = x^{n+1} y^{n+1} [y, x]^{\frac{(n+1)n}{2}}. \qh
    \end{align*}
\end{sol}

\newpage

\begin{exercise}[5.4.9]
    Prove that if $p$ is an odd prime, and $P$ is a group of order $p^3$, then the $p$-th power map $x \mapsto x^p$ is a homomorphism of $P$ into $Z(P)$. If $P$ is not cyclic, show that the kernel of the $p$-th power map has order $p^2$ or $p^3$. Is the squaring map a homomorphism in non-abelian groups of order 8? Where is the oddness of $p$ needed in the above proof?
\end{exercise}

\begin{sol}
    Note that the abelian case is trivial, so we may assume $P$ is non-abelian. By \exref{5.4.7}, we have $P' = Z(P)$, so $P/Z(P)$ is abelian. Note that $[y, x] \in P' = Z(P)$ for all $x, y \in P$, so $x$ and $y$ both commute with $[y, x]$. Moreover, $|Z(P)| = p$ by \exref{5.4.7}, hence
    \[(xy)^p = x^p y^p [y, x]^{\frac{p(p-1)}{2}} = x^p y^p,\]
    where the last equality holds since $[y, x]$ has order dividing $p$ and $\frac{p(p-1)}{2}$ is divisible by $p$ when $p$ is odd. Thus, the $p$-th power map is a homomorphism of $P$ into $Z(P)$. To see that $x^p \in Z(P)$ for all $x \in P$, note that for any $y \in P$, we have
    \[y x^py\inv = (y x y\inv)^p = (x [y, x])^p = x^p [y, x]^p = x^p.\]
    Suppose now that $P$ is not cyclic. Since the map is to $Z(P)$, which has order $p$, the image of the map has order dividing $p$. If the image is trivial, then the kernel is all of $P$, so the kernel has order $p^3$. If the image is non-trivial, then the kernel has order $p^2$ by the First Isomorphism Theorem.

    For the squaring map in non-abelian groups of order 8, we have $Q_8$ and $D_8$. In $Q_8$, observe that $(ij)^2 = (k)^2 = -1$, while $i^2j^2 = 1$. In $D_8$, we similarly have $(rs)^2 = 1$, while $r^2s^2 = r^2$. Thus, the squaring map is not a homomorphism in non-abelian groups of order 8. Lastly, the oddness of $p$ is needed in the proof to ensure that $\frac{p(p-1)}{2}$ is divisible by $p$, which allows us to conclude that $[y, x]^{\frac{p(p-1)}{2}} = 1$.
\end{sol}

\begin{specialexercise}[5.4.10]
    Prove that a finite abelian group is the direct product of its Sylow subgroups.
\end{specialexercise}

\begin{sol}
    Inducting on the number of prime factors of the group, we first consider $n = 1$, or when $|G| = p_1^n$ for some prime $p_1$. Then $G$ is its own Sylow $p_1$-subgroup, so the statement holds. We now assume that any abelian with fewer than $n$ prime factors is the direct product of its Sylow subgroups, and let $G$ be an abelian group with $n$ prime factors. Put
    \[|G| = p_1^{\alpha_1} p_2^{\alpha_2} \cdots p_n^{\alpha_n}.\]
    Let $H$ be some subgroup of $G$ generated by the Sylow $p_i$ subgroups for $1 \leq i \leq n - 1$. By definition, every element of $H$ is generated by some element from a Sylow $p_i$-subgroup, so its order must divide the lcm $p_1^{\alpha_1} p_2^{\alpha_2} \cdots p_{n-1}^{\alpha_{n-1}}$, so $H$ is a proper subgroup of $G$. Moreover, $H$ contains all such Sylow $p_i$-subgroups for $1 \leq i \leq n - 1$, hence $|H| \geq p_1^{\alpha_1} p_2^{\alpha_2} \cdots p_{n-1}^{\alpha_{n-1}}$. Since $H$ is a proper subgroup of $G$, we must have $|H| = p_1^{\alpha_1} p_2^{\alpha_2} \cdots p_{n-1}^{\alpha_{n-1}}$. By the inductive hypothesis, $H$ is the direct product of its Sylow subgroups, so
    \[H = P_1 \times P_2 \times \cdots \times P_{n-1},\]
    where $P_i$ is the Sylow $p_i$-subgroup of $G$ for $1 \leq i \leq n - 1$. Since $G$ is abelian, we have that $P_n$ is normal in $G$, so $HP_n$ is a subgroup of $G$. Moreover, since $H$ and $P_n$ have trivial intersection because $p_n^{\alpha_n}$ and $p_1^{\alpha_1} p_2^{\alpha_2} \cdots p_{n-1}^{\alpha_{n-1}}$ are coprime, we have that $HP_n = H \times P_n$. Since $|HP_n| = |H||P_n| = |G|$, we must have $HP_n = G$, so
    \[G = P_1 \times P_2 \times \cdots \times P_n.\] \qh
\end{sol}

\newpage

\begin{specialexercise}[5.4.11]
    Prove that if $G = HK$ where $H$ and $K$ are characteristic subgroups of $G$ with $H \cap K = 1$, then $\aut(G) \cong \aut(H) \times \aut(K)$. Deduce that if $G$ is an abelian group of finite order, then $\aut(G)$ is isomorphic to the direct product of the automorphism groups of its Sylow subgroups.
\end{specialexercise}

\begin{sol}
    Let $\phi \in \aut(G)$. Since $H$ and $K$ are characteristic subgroups of $G$, we have $\phi(H) = H$ and $\phi(K) = K$. Thus, we can define $\phi_H = \phi|_H$ and $\phi_K = \phi|_K$, which are automorphisms of $H$ and $K$ respectively. We can then define a map $\Phi: \aut(G) \to \aut(H) \times \aut(K)$ by $\Phi(\phi) = (\phi_H, \phi_K)$. To see that $\Phi$ is a homomorphism, let $\phi, \psi \in \aut(G)$ be given. Then
    \[\Phi(\phi \circ \psi) = ((\phi \circ \psi)_H, (\phi \circ \psi)_K) = (\phi_H \circ \psi_H, \phi_K \circ \psi_K) = \Phi(\phi) \circ \Phi(\psi).\]
    To see that $\Phi$ is injective, let $\phi \in \aut(G)$ be given such that $\Phi(\phi) = (1_H, 1_K)$, where $1_H$ and $1_K$ are the identity automorphisms of $H$ and $K$ respectively. Then $\phi(h) = h$ for all $h \in H$ and $\phi(k) = k$ for all $k \in K$. Since every element of $G$ can be uniquely expressed as $hk$ for some $h \in H$ and $k \in K$, we have $\phi(hk) = \phi(h)\phi(k) = hk$, so $\phi$ is the identity automorphism of $G$. Thus, $\Phi$ is injective.

    To show that $\Phi$ is surjective, let $(\psi_H, \psi_K) \in \aut(H) \times \aut(K)$ be given. We can define a map $\psi: G \to G$ by $\psi(hk) = \psi_H(h)\psi_K(k)$ for all $h \in H$ and $k \in K$. Because $H \cap K = 1$, and $G = HK$, every element of $G$ can be uniquely expressed as $hk$ for some $h \in H$ and $k \in K$. We check that $\psi$ is a homomorphism:
    \[\psi(h_1k_1 \cdot h_2k_2) = \psi(h_1h_2k_1k_2) = \psi_H(h_1h_2)\psi_K(k_1k_2) = \psi_H(h_1)\psi_H(h_2)\psi_K(k_1)\psi_K(k_2) = \psi(h_1k_1)\psi(h_2k_2).\]
    where $k_1$ and $h_2$ commute since $H$ and $K$ are the factors of $G = H \times K$. Since $\psi_H$ and $\psi_K$ are automorphisms, $\psi$ is a bijection, hence an automorphism of $G$. Moreover, $\Phi(\psi) = (\psi_H, \psi_K)$, so $\Phi$ is surjective.

    Therefore, we have shown that $\Phi: \aut(G) \to \aut(H) \times \aut(K)$ is an isomorphism. For the second part of the problem, since a finite abelian group is the direct product of its Sylow subgroups by the previous exercise, and conjugation is trivial in an abelian group, we have that each Sylow subgroup is characteristic, so the automorphism group of the finite abelian group is isomorphic to the direct product of the automorphism groups of its Sylow subgroups.
\end{sol}

\begin{exercise}[5.4.12]
    Use Theorem 4.17 to describe the automorphism group of a finite cyclic group.
\end{exercise}

\begin{sol}
    From Proposition 4.17, we know that $\aut(Z_p) \cong Z_{p - 1}$, and $\aut(Z_{p^n}) \cong Z_{p^{n - 1}(p - 1)}$ for any prime $p$ and $n \geq 2$. For $n \geq 3$ and $p = 2$, we have that $\aut(Z_{2^n}) \cong Z_2 \times Z_{2^{n - 2}}$ that is not cyclic. Hence, for any finite cyclic group $Z_n$ with (specific) prime factorization $n = 2^{\alpha_1} p_2^{\alpha_2} \cdots p_k^{\alpha_k}$, we have by \exref{5.4.10} and the previous exercise that
    \[\aut(Z_n) \cong \aut(Z_{2^{\alpha_1}}) \times \aut(Z_{p_2^{\alpha_2}}) \times \cdots \times \aut(Z_{p_k^{\alpha_k}}).\]
    Note that $\aut(Z_{p_i^{\alpha_i}}) \cong Z_{p_i^{\alpha_i - 1}(p_i - 1)}$ is cyclic for $i \geq 2$, and $\aut(Z_{2^{\alpha_1}})$ is cyclic if $\alpha_1 \leq 2$ and is not cyclic if $\alpha_1 \geq 3$. Thus, the structure of $\aut(Z_n)$ depends on the exponent of 2 in the prime factorization of $n$. We can check the following cases:
    \begin{itemize}
        \item $\alpha_1 = 0$: Then the 2-part of $n$ is trivial, so $\aut(Z_n) \cong Z_{p_2^{\alpha_2 - 1}(p_2 - 1)} \times \cdots \times Z_{p_k^{\alpha_k - 1}(p_k - 1)}$ is cyclic if and only if $k = 1$, or the $p_i^{\alpha_i - 1}(p_i - 1)$ are pairwise coprime.
        \item $\alpha_1 = 1$: Then the 2-part of $n$ is $Z_2$ with $\aut(Z_2)$ trivial, so $\aut(Z_n) \cong Z_{p_2^{\alpha_2 - 1}(p_2 - 1)} \times \cdots \times Z_{p_k^{\alpha_k - 1}(p_k - 1)}$ is cyclic if and only if the orders $p_i^{\alpha_i - 1}(p_i - 1)$ are pairwise coprime and coprime to 2, which is equivalent to the condition in the previous case. Hence, the same condition as $\alpha_1 = 0$ applies here.
        \item $\alpha_1 = 2$: Then the 2-part of $n$ is $Z_4$, so $\aut(Z_n) \cong \aut(Z_4) \times Z_{p_2^{\alpha_2 - 1}(p_2 - 1)} \times \cdots \times Z_{p_k^{\alpha_k - 1}(p_k - 1)}$ is cyclic if and only if $k = 1$, similar to $\alpha_1 = 0$, or the $p_i^{\alpha_i - 1}(p_i - 1)$ are pairwise coprime and coprime to 2.
        \item $\alpha_1 \geq 3$: Then the 2-part of $n$ is $Z_2 \times Z_{2^{\alpha_1 - 2}}$, so $\aut(Z_n) \cong \aut(Z_{2^{\alpha_1}}) \times Z_{p_2^{\alpha_2 - 1}(p_2 - 1)} \times \cdots \times Z_{p_k^{\alpha_k - 1}(p_k - 1)}$ is not cyclic. \qh
    \end{itemize}
\end{sol}

\begin{exercise}[5.4.13]
    Prove that $D_{8n}$ is not isomorphic to $D_{4n} \times Z_2$ for any $n \geq 1$.
\end{exercise}

\begin{sol}
    Note that $Z(D_{8n}) = \set{1, r^{2n}}$ has order 2, while $Z(D_{4n} \times Z_2) = Z(D_{4n}) \times Z(Z_2)$ has order 4, so $D_{8n}$ cannot be isomorphic to $D_{4n} \times Z_2$ for any $n \geq 1$.
\end{sol}

\begin{exercise}[5.4.14]
    Let $G = \set{(a_{ij}) \in \gl_n(F) \mid a_{ij} = 0 \text{ for } i > j, \text{ and } a_{11} = a_{22} = \cdots = a_{nn}}$, where $F$ is a field, be the group of upper triangular matrices all of whose diagonal entries are equal. Prove that $G \cong D \times U$, where $D$ is the group of all nonzero multiples of the identity matrix, and $U$ is the group of upper triangular matrices with 1's down the diagonal.
\end{exercise}

\begin{sol}
    Note that $I_n \in D$. For any $\lambda, \mu \in F\unt$, then $(\lambda I_n)(\mu I_n)\inv = (\lambda I_n)(\mu\inv I_n) = (\lambda\mu\inv)I_n \in D$, so $D$ is a subgroup of $G$. For $A, B \in U$, we know that the product still has $AB$ upper triangular with 1's down the diagonal, so $U$ is a subgroup of $G$. To see that $D \nsub G$, let $A \in G$ be given. Then $A$ is upper triangular with all diagonal entries equal to some $\lambda \in F\unt$. For any $\mu I_n \in D$, we have
    \[A(\mu I_n)A\inv = \mu AI_n A\inv = \mu A A\inv = \mu I_n \in D,\]
    so $D$ is normal in $G$. To see that $U \nsub G$, let $A \in G$ be given. Then $A$ is upper triangular with all diagonal entries equal to some $\lambda \in F\unt$. We may factor out $\lambda = a_{11}$ from $A$ to get $A = a_{11} X$, where $X$ is upper triangular with 1's down the diagonal, so $X \in U$. For any $Y \in U$, we have
    \[A Y A\inv = a_{11} X Y X\inv a_{11}\inv = X Y X\inv,\]
    where $XYX \inv \in U$ since $U \leq G$, so $U$ is normal in $G$. If $Z \in D \cap U$, then $Z = \lambda I_n$ for some $\lambda \in F\unt$, and $Z$ is upper triangular with 1's down the diagonal, so $\lambda = 1$, hence $Z = I_n$. Thus, $D \cap U = 1$. Lastly, for any $A \in G$, we can factor out the common diagonal entry to get $A = a_{11} X$ for some $X \in U$, so $G = DU \cong D \times U$ by Theorem 5.9.
\end{sol}

\begin{exercise}[5.4.15]
    If $A$ and $B$ are normal subgroups of $G$ such that $G/A$ and $G/B$ are both abelian, prove that $G/(A \cap B)$ is abelian.
\end{exercise}

\begin{sol}
    $G/A$ abelian implies $G' \leq A$, and $G/B$ abelian implies $G' \leq B$, so $G' \leq A \cap B$. Hence, $G/(A \cap B)$ is abelian by Proposition 5.7(4).
\end{sol}

\begin{exercise}[5.4.16]
    Prove that if $K$ is a normal subgroup of $G$, then $K' \nsub G$.
\end{exercise}

\begin{sol}
    Let $K$ be a normal subgroup of $G$. By Proposition 5.7(2), we have $[G, K] \leq K$. Since $K' = [K, K] \leq [G, K]$, we have $K' \leq K$. Moreover, for any $g \in G$ and $x = c_1 c_2 \cdots c_n \in K'$, where $c_i = [k_i, \ell_i]$ for some $k_i, \ell_i \in K$, we have
    \[g x g\inv = g c_1 c_2 \cdots c_n g\inv = (g c_1 g\inv)(g c_2 g\inv) \cdots (g c_n g\inv),\]
    where $g c_i g\inv = g [k_i, \ell_i] g\inv = [g k_i g\inv, g \ell_i g\inv]$ since $K$ is normal in $G$, so $g x g\inv \in K'$. Thus, $K' \nsub G$.
\end{sol}

\begin{exercise}[5.4.17]
    If $K$ is a normal subgroup of $G$ and $K$ is cyclic, prove that $G' \leq C_G(K)$. [Recall that the automorphism group of a cyclic group is abelian.]
\end{exercise}

\begin{sol}
    Let $K$ be a normal subgroup of $G$ and $K$ be cyclic. Since $K$ is cyclic, we have $\aut(K)$ is abelian. Moreover, the homomorphism of $G$ into $\aut(K)$ given by conjugation has kernel $C_G(K)$, so $G/C_G(K)$ is isomorphic to a subgroup of $\aut(K)$, which is abelian. Hence, $G' \leq C_G(K)$ by Proposition 5.7(4).
\end{sol}

\begin{exercise}[5.4.18]
    Let $K_1, K_2, \ldots, K_n$ be non-abelian simple groups, and let $G = K_1 \times K_2 \times \cdots \times K_n$. Prove that every normal subgroup of $G$ is of the form $G_I$ for some subset $I$ of $\set{1, 2, \ldots, n}$ (where $G_I$ is defined in \exref{5.1.2}). [If $N \nsub G$ and $x = (a_1, \ldots, a_n) \in N$ with some $a_i \neq 1$, then show that there is some $g_i \in G_i$ not commuting with $a_i$. Show that $[(1, \ldots, g_i, \ldots, 1), x] \in K_i \cap N$, and deduce $K_i \leq N$.]
\end{exercise}

\begin{sol}
    Let $N$ be a normal subgroup of $G$. If $N = 1$, then $N = G_\emptyset$. If $N = G$, then $N = G_{\set{1, 2, \ldots, n}}$. Now, assume that $1 < N < G$. Then there exists some $x = (a_1, a_2, \ldots, a_n) \in N$ such that $a_i \neq 1$ for some $i$. Since $K_i$ is non-abelian simple, there exists some $g_i \in K_i$ such that $g_i$ does not commute with $a_i$. Consider the element $(1, \ldots, g_i, \ldots, 1) \in G$, where $g_i$ is in the $i$-th position. We have
    \[(1, \ldots, g_i, \ldots, 1) x (1, \ldots, g_i\inv, \ldots, 1) = (a_1, \ldots, g_i a_i g_i\inv, \ldots, a_n).\]
    Hence,
    \[(1, \ldots, g_i, \ldots, 1) x (1, \ldots, g_i\inv, \ldots, 1) x\inv = (1, \ldots, [g_i, a_i], \ldots, 1) \in N.\]
    Since $[g_i, a_i]$ is non-trivial and is an element of the simple group $K_i$, we have that $\langle [g_i, a_i]^{K_i} \rangle \nsub K_i$. However, simplicity of $K_i$ forces $\langle [g_i, a_i]^{K_i} \rangle = K_i$, so $K_i \leq N$. By repeating this argument for any other non-trivial component of an element of $N$, we can find some subset $I$ of $\set{1, 2, \ldots n}$ such that $G_I = K_{i_1} \times K_{i_2} \times \cdots K_{i_k} \leq N$, where $\set{i_1, i_2, \ldots i_k} = I$, which shows $G_I \leq N$. On the other hand, if $N$ contains some element $x = (a_1, a_2, \ldots, a_n)$ such that $a_j \neq 1$ for some $j \notin I$, then we can similarly find some element of $N$ that has non-trivial component in the $j$-th position, which contradicts that $K_j \not\leq N$. Thus, $N \leq G_I$. Therefore, $N = G_I$.
\end{sol}

\begin{exercise}[5.4.19]
    A group $H$ is called \emph{perfect} if $H' = H$ (i.e., $H$ equals its own commutator subgroup).
    \begin{subexercises}
        \item Prove that every non-abelian simple group is perfect.
        \item Prove that if $H$ and $K$ are perfect subgroups of a group $G$, then $\gen{H, K}$ is also perfect. Extend this to show that the subgroup of $G$ generated by any collection of perfect subgroups is perfect.
        \item Prove that any conjugate of a perfect subgroup is perfect.
        \item Prove that any group $G$ has a maximal perfect subgroup, and that this subgroup is normal.
    \end{subexercises}
\end{exercise}

\begin{solenum}
    \item Let $H$ be a non-abelian simple group. Since $H$ is non-abelian, we have $H' \neq 1$. Since $H$ is simple, the only normal subgroups of $H$ are $1$ and $H$, so we must have $H' = H$.
    \item Let $H$ and $K$ be perfect subgroups of a group $G$. We want to show that $\gen{H, K}' = \gen{H, K}$. Since $H$ and $K$ are perfect, we have $H' = H$ and $K' = K$, so $H \leq H'$ and $K \leq K'$. Moreover, since $H$ and $K$ are subgroups of $\gen{H, K}$, we have $H' \leq \gen{H, K}'$ and $K' \leq \gen{H, K}'$, so $H \leq \gen{H, K}'$ and $K \leq \gen{H, K}'$. Hence, $\gen{H, K} = \gen{H, K}'$. The extension to any collection of perfect subgroups follows by induction on the number of perfect subgroups in the collection.
    \item Let $H$ be a perfect subgroup of a group $G$, and let $g \in G$ be given. We want to show that $g\inv H g$ is perfect. Since $H$ is perfect, we have $H' = H$, so $g\inv H' g = g\inv H g$. Moreover, since conjugation is an automorphism of $G$, we have $g\inv H' g = (g\inv H g)'$, so $(g\inv H g)' = g\inv H g$. Hence, $g\inv H g$ is perfect.
    \item Let $\mathcal{P}$ be the collection of perfect subgroups of $G$. Since $1$ is a perfect subgroup of $G$, we have $\mathcal{P}$ is nonempty. Let $\mathcal{C}$ be a chain in $\mathcal{P}$, and let $H = \bigcup_{K \in \mathcal{C}} K$. To see that $H \leq G$, note that for any $x, y \in H$, then there exists $K_x, K_y \in \mathcal C$ such that $x \in K_x$ and $y \in K_y$. Without loss of generality, we may assume $K_x \leq K_y$, so $x, y \in K_y$, hence $xy\inv \in K_y \leq H$. Thus, $H$ is a subgroup of $G$. To see that $H \in \pp$, observe that for any $x \in H$, then $x \in K$ for some $K \in \mathcal{C}$, so $x \in K' = K$, hence $x \in H'$. Thus, $H \leq H'$. Since $H' \leq H$ trivially, then $H = H'$, hence $H \in \mathcal P$. By Zorn's Lemma, $\mathcal{P}$ has a maximal element $M$. 
    
    To see that $M$ is normal in $G$, let $g \in G$ be given. Then $g\inv M g$ is perfect by part (c), so $g\inv M g \in \mathcal{P}$. Since $M$ is perfect, we may use part (b) to get that $\gen{M, g\inv M g}$ is perfect, so $\gen{M, g\inv M g} \in \mathcal{P}$. Since $M \leq \gen{M, g\inv M g}$ and $M$ is maximal in $\mathcal{P}$, we must have $M = \gen{M, g\inv M g}$, so $g\inv M g \leq M$. By symmetry of the argument, we also have $M \leq g\inv M g$, so $g\inv M g = M$. Hence, $M$ is normal in $G$.
\end{solenum}

\begin{exercise}[5.4.20]
    Let $H(F)$ be the Heisenberg group over the field $F$, cf. \exref{1.4.11}. Find an explicit formula for the commutator $[X, Y]$, where $X, Y \in H(F)$, and show that the commutator subgroup of $H(F)$ equals the center of $H(F)$ (cf. \exref{2.2.14}).
\end{exercise}

\begin{sol}
    Let
    \[X = 
    \begin{bmatrix}
        1 & a & b \\
        0 & 1 & c \\
        0 & 0 & 1
    \end{bmatrix}, \quad
    Y = 
    \begin{bmatrix}
        1 & d & e \\
        0 & 1 & f \\
        0 & 0 & 1
    \end{bmatrix} \in H(F).\]
    Then we have
    \[X\inv =
    \begin{bmatrix}
        1 & -a & ac - b \\
        0 & 1 & -c \\
        0 & 0 & 1
    \end{bmatrix}, \quad
    Y\inv =
    \begin{bmatrix}
        1 & -d & df - e \\
        0 & 1 & -f \\
        0 & 0 & 1
    \end{bmatrix}.\]
    Hence,
    \begin{align*}
        [X, Y] & = X\inv Y\inv XY \\
        & =
        \begin{bmatrix}
            1 & -a & ac - b \\
            0 & 1 & -c \\
            0 & 0 & 1
        \end{bmatrix}
        \begin{bmatrix}
            1 & -d & df - e \\
            0 & 1 & -f \\
            0 & 0 & 1
        \end{bmatrix}
        \begin{bmatrix}
            1 & a & b \\
            0 & 1 & c \\
            0 & 0 & 1
        \end{bmatrix}
        \begin{bmatrix}
            1 & d & e \\
            0 & 1 & f \\
            0 & 0 & 1
        \end{bmatrix} \\
        & =
        \begin{bmatrix}
            1 & 0 & af - cd \\
            0 & 1 & 0 \\
            0 & 0 & 1
        \end{bmatrix}
    \end{align*}
    Recall that the center of $H(F)$ consists of all matrices of the form
    \[
    \begin{bmatrix}
        1 & 0 & z \\
        0 & 1 & 0 \\
        0 & 0 & 1
    \end{bmatrix}\]
    for some $z \in F$. Clearly, $af - cd \in F$, so $[X, Y] \in Z(H(F))$, hence $H(F)' \leq Z(H(F))$. On the other hand, for any $z \in F$, we may choose $a = 1$, $f = z$, and $c = d = 0$ to get that $Z(H(F)) \leq H(F)'$. Therefore, $H(F)' = Z(H(F))$.
\end{sol}